Un an après, c'était la paix d'Alès. En ce moment-là, notre Pujol
de la Taie\footnote{de la Tapie?} est sur le point de s'emparer de Viane qui a tenu en
échec Condé. Le 9 août 1629, Richelieu arrive de Pézenas. Il y a
dissous le parlement du Languedoc. Le régionalisme cause quelques
ennuis. L'éminent ministre de Louis XIII traverse notre petite
ville. Le même jour, dans la soirée, le Cardinal de Richelieu
fait son entrée à Albi.

Les épisodes comiques ou tragiques se sont succédés au cours de
cette guerre civile qui ensanglanta nos murs. Dans un cahier
détenu par les descendants de la famille Pujol de la Tapie, on
peut lire le trait suivant : Jean Pujol était venu pour
« pétarder » une des maisons de Peyreségade où se sont réunis
des soldats protestants de Viane. Au moment où il se prépare à
mettre à exécution son projet, il entend les militaires ennemis
boire à sa santé, à la santé de « l'artillogat ». Un des buveurs
paraît même à la fenêtre et lui « pisse » sur le chapeau. Au
milieu des rires, on entend l'auteur de cet exploit s'écrier :
« je viens de pisser sur le chapeau d'une personne ». Notre
Pujol se croit trahi et détale... la mèche de son pétard
était-elle quelque peu humide à la suite de cet incident ?...
On peut le penser !!!